\section{Related work}
\label{sec:related}

\paragraph{Open-loop evaluation.}
Codevilla~\etal~\cite{codevilla2018offline} showed that offline metrics of vision-based driving
models correlate weakly with driving quality.
For end-to-end planning on nuScenes, an MLP on ego status alone matches camera-based planners on
displacement metrics~\cite{zhai2023admlp,li2024egostatus}, which shows that displacement can be
optimised without perception.

\paragraph{Pseudo-simulation.}
NAVSIM~\cite{dauner2024navsim} scores a planned trajectory by simulating it for four seconds
against a non-reactive world with the PDM Score~\cite{dauner2023pdm} built on nuPlan~\cite{caesar2021nuplan}.
Its second version~\cite{cao2025pseudosim} adds a second stage of synthetic scenes near where the
first plan leads and reports high agreement with closed-loop results.
On non-reactive navtest, even the choice of waypoint controller is rewarded for fitting the fixed
future~\cite{rodriguez2026waypoint}.

\paragraph{Closed-loop evaluation.}
Bench2Drive~\cite{jia2024bench2drive} evaluates CARLA-trained policies over fixed routes; HUGSIM~\cite{zhou2024hugsim}
replays reconstructed real scenes; BridgeSim~\cite{zhao2026bridgesim} converts NAVSIM scenes into
MetaDrive~\cite{li2022metadrive} and runs NAVSIM-trained checkpoints in closed loop.
BridgeSim documents an open-loop/closed-loop gap for policies; we use it as a reference to test
evaluators.

\paragraph{Evaluating evaluators.}
Wang~\etal~\cite{wang2026olcl} correlate NAVSIM and Bench2Drive scores of eight published methods.
In robot manipulation, SIMPLER~\cite{li2024simpler} measures whether simulated scores preserve real
policy rankings, introducing the mean maximum rank violation we also report.
None of these has a known ordering of policies; ours does, by construction.
